\documentclass[11pt]{article}
\usepackage[margin=0.85in]{geometry}
\usepackage[T1]{fontenc}
\usepackage{lmodern}
\usepackage{microtype}
\usepackage{booktabs}
\usepackage{longtable}
\usepackage{array}
\usepackage{enumitem}
\usepackage{xcolor}
\usepackage{hyperref}
\usepackage{url}
\hypersetup{colorlinks=true,linkcolor=blue!55!black,citecolor=blue!55!black,urlcolor=blue!60!black}
\setlist{nosep,leftmargin=*}
\newcommand{\alem}{\textsc{Alem}}
\newcommand{\dice}{\textsc{AlemDICE}}
\newcommand{\must}{\textbf{Must}}
\newcommand{\should}{\textbf{Should}}
\newcommand{\could}{\textbf{Could}}

\title{Improving \dice{}: A Research Review and Experimental Roadmap}
\author{Repository technical report}
\date{21 July 2026}

\begin{document}
\maketitle

\begin{abstract}
\alem{} shows that frontier language models can make base-task progress while still failing at
coordination, especially when actions must be synchronized.  This report reviews the published
\alem{} study, related work in LLM-agent coordination, cooperative AI, long-horizon agent design,
evaluation methodology, and LLM serving.  It then audits the current \dice{} harness and proposes a
prioritized improvement program.  The highest-value near-term change is not a larger prompt or a
central planner.  It is an explicit, decentralized commitment protocol backed by structured plan
memory, event-triggered messages, deterministic action checks, and paired-seed evaluation.  The
report also recommends partner-generalization tests, crash-stop experiments, and a serving study
that separates policy quality from infrastructure failures.  Every proposal preserves an unchanged
published baseline and is framed as a falsifiable ablation rather than an assumed improvement.
\end{abstract}

\section{Scope and method}

The object of improvement is the complete text-agent system: observation and action interface,
reasoning policy, communication, memory, execution scheduling, evaluation, and local inference.
The environment's scientific value depends on keeping the canonical three-agent \alem{} control
runnable.  Consequently, recommendations in this report are new experiment profiles; none should
silently change the leaderboard protocol.

The review uses four evidence classes.  First, it reads the complete 42-page \alem{} preprint and its
appendices \cite{alem2026}.  Second, it examines the current repository configuration, evaluator,
agent parser, communication tracker, team-leader treatment, provenance ledger, and Ollama provider.
Third, it compares findings with primary work on coordination benchmarks, partner generalization,
long-horizon agents, and statistical evaluation
\cite{agashe2025coordination,zhu2025multiagentbench,sun2025collab,leibo2021meltingpot,agarwal2021precipice}.
Fourth, it considers serving evidence and the observed four-A100 deployment
\cite{kwon2023vllm,ollamafaq,ollamacontext}.  Claims explicitly distinguish published results,
repository facts, and proposed hypotheses.

\section{What the evidence says}

\subsection{The benchmark exposes a real coordination gap}

\alem{} combines procedural worlds with coordination requirements spanning delayed dependencies,
handovers, and same-timestep joint action.  Thirteen homogeneous LLM teams average only about 6\%
normalized return.  The strongest base-task model is not necessarily the strongest coordinator:
GPT-5.4-High has strong base reward but substantially weaker coordination reward, while
Gemini-3.1-Pro-High approaches a one-billion-step MARL reference on Hard coordination
\cite{alem2026}.  This supports reporting base and coordination outcomes separately.

The most diagnostic failures are synchronous-hard tasks.  Success requires discovering a common
target, navigating to compatible positions, communicating roles and timing, and issuing compatible
actions on the same tick.  Even the best evaluated LLM covers only a minority of such rewards.
Handover tasks are easier because their time window tolerates imperfect agreement.  A productive
harness should therefore target agreement and execution timing, not merely improve navigation.

\subsection{Communication helps, but free-form broadcast is inefficient}

Removing communication produces the largest coordination drop in the paper: Gemini coordination
falls from 17.5\% to 5.3\%, and Gemma from 8.8\% to 3.8\% on Hard \cite{alem2026}.  Yet the
qualitative analysis shows that messages are mostly intent announcements and commands; questions
and acknowledgements are rare.  Gemma messages often behave as local status reports.  Free-form
broadcast therefore transmits useful information without guaranteeing that recipients formed the
same commitment.

This is consistent with LLM-Coordination, which finds weaknesses when success depends on modeling
partners' beliefs and joint plans rather than environmental facts \cite{agashe2025coordination}.
MultiAgentBench finds that topology and cognitive planning affect milestone completion
\cite{zhu2025multiagentbench}; Collab-Overcooked similarly argues for process metrics because goal
interpretation can look strong while active collaboration and adaptation remain weak
\cite{sun2025collab}.  Together these results favor measuring delivered coordination state, not raw
message volume.

\subsection{Memory helps only when it preserves plans}

The \alem{} scratchpad ablation is model-dependent.  Gemini's entries frequently preserve
future-oriented, multi-line, sometimes turn-indexed plans, and removing memory hurts its
coordination.  Gemma's entries are usually one-line coordinate-heavy state summaries; removing
them has little measured effect.  Disabling Gemma reasoning makes its scratchpad longer and more
deliberative, but does not recover performance \cite{alem2026}.  More tokens are therefore not the
right objective.  Memory should preserve commitments, unresolved requests, discovered facts, and
failure evidence while avoiding facts already present in the observation.

ReAct supports tight interleaving of reasoning and grounded action \cite{yao2023react}, while
Voyager demonstrates the utility of persistent, executable skill abstractions in an open-ended world
\cite{wang2024voyager}.  For \dice{}, the conservative lesson is to build a small typed plan ledger
and a tested skill layer before attempting unconstrained cross-episode reflection.

\subsection{Heterogeneity and centralization are not free wins}

The paper's mixed-model teams perform near the average of their constituent homogeneous teams,
not the strongest member \cite{alem2026}.  Different planning horizons and message conventions can
create new coordination costs.  A bodyless leader can make assignments, but it also introduces a
privileged aggregation point, serial latency, extra tokens, and a single point of failure.  The current
\dice{} leader treatment is scientifically useful as a comparison; it should not replace the
decentralized headline.

Partner diversity remains important.  Overcooked studies zero-shot coordination with new partners
\cite{carroll2019overcooked}, and Melting Pot makes generalization to new co-player policies a first-
class evaluation goal \cite{leibo2021meltingpot}.  A robust \dice{} policy must work with unfamiliar
message styles and competence levels, not only three copies of itself.

\section{Audit of the current \dice{} baseline}

\begin{longtable}{p{0.20\linewidth}p{0.35\linewidth}p{0.37\linewidth}}
\toprule
Area & Current strength & Limitation or confound \\
\midrule
\endhead
Environment & Procedural coordination, soft specialization, three difficulty levels, matched seeds. & Most LLM runs die before advanced tiers; aggregate score can hide opportunity exposure. \\
Observation & Compact text, affordances, progressive disclosure, local history. & Spatial state remains prose-heavy; no explicit uncertainty or change set; possible leakage needs systematic audit. \\
Action & Canonical action list, tagged parser, fallback to Noop, parse metrics. & Syntactic validity is checked after generation; semantic feasibility and repeated-failure recovery are weak. \\
Communication & Parallel worker calls, delayed peer broadcast, byte and lexical metrics. & Broadcast is unstructured; no recipient, message type, commitment ID, expiry, acknowledgement, or conflict resolution. \\
Memory & Private bounded scratchpad and reasoning history. & Free text encourages state restatement; no durable goal, commitment, evidence, or confidence fields. \\
Planning & Model-native reasoning; optional bodyless leader. & Baseline has no explicit decentralized plan agreement; leader serializes a phase and changes information topology. \\
Evaluation & Durable attempt ledger, provenance, termination reasons, usage and latency accounting. & Twenty worlds still yield wide uncertainty; opportunity-conditioned and partner-generalization metrics are incomplete. \\
Serving & Native Ollama client, retries, separated reasoning, one endpoint per agent possible. & Readiness checks do not guarantee model load; auto context and parallel slots can cause VRAM failures; provider retries can mask policy latency. \\
\bottomrule
\end{longtable}

The implementation already contains important safeguards: per-agent clients, parallel worker
futures, immutable configuration snapshots, failed-attempt accounting, and communication routing
metrics.  Improvements should build on those observability hooks instead of adding a second opaque
agent loop.

\section{Prioritized improvements}

\subsection{P0: a decentralized commitment protocol}

The highest-priority experiment should replace unconstrained coordination prose with an optional
typed envelope while retaining a short free-text payload:

\begin{verbatim}
TYPE: PROPOSE | ACCEPT | REJECT | READY | COMMIT | CANCEL | STATUS
ID: stable task/commitment identifier
TO: agent ids or team
TARGET: entity and coordinate
ACTION: canonical action
WINDOW: earliest/latest tick
PRECONDITIONS: local facts
PAYLOAD: <= 160 characters
\end{verbatim}

A synchronous task uses a two-phase handshake: one agent proposes a target and tick window;
partners explicitly accept or reject; all send READY after reaching valid positions; then each commits
to the same action/tick.  Conflicting commitments are surfaced in the next prompt.  Expired
commitments are removed automatically.  This mechanism remains decentralized: no process chooses
the plan, and each message contains only sender-visible information.

\textbf{Hypothesis.} Explicit acknowledgement and stable IDs will improve synchronous-hard
opportunity conversion and reduce repeated contradictory messages without increasing total bytes.
\textbf{Primary metrics.} Proposal acceptance latency, agreement rate, ready-to-execution latency,
commitment violations, synchronized-action success per exposed opportunity, delivered bytes, and
Coord.\%.  \textbf{Control.} The unchanged free-form broadcast baseline with identical model,
prompts outside the protocol block, seeds, and generation parameters.

\subsection{P0: structured private plan memory}

Replace the optional free-text scratchpad in one ablation with a bounded schema:

\begin{verbatim}
GOAL: current role-consistent objective
NEXT: up to three state-contingent steps
COMMITMENTS: ids, counterparties, target, tick/window, status
FACTS: off-screen discoveries with source and last-confirmed tick
FAILURES: attempted action, observed outcome, revised belief
\end{verbatim}

The harness should reject duplicate observation text, age facts, and cap each section separately.
Memory updates should be patches rather than full rewrites, reducing both drift and token cost.  This
directly tests the paper's inference that planner-like memory is useful while state restatement is not.
Measure stale-fact use, commitment recall, plan completion, memory novelty, and tokens in addition
to reward.

\subsection{P0: opportunity-conditioned diagnostics}

Reward coverage alone cannot distinguish ``the team never found the site'' from ``the team found it
and failed to coordinate.''  Instrument a coordination opportunity state machine:

\begin{enumerate}
  \item opportunity exists in world;
  \item at least one agent observes it;
  \item information reaches required partners;
  \item a compatible plan is proposed;
  \item agents arrive in valid positions;
  \item agreement is reached;
  \item joint action is attempted;
  \item environment confirms success.
\end{enumerate}

Report transition probabilities and time in each state by task class.  These metrics localize whether
an intervention improves discovery, routing, navigation, agreement, or execution.  They also align
with process-oriented evaluation in Collab-Overcooked and milestone metrics in MultiAgentBench
\cite{sun2025collab,zhu2025multiagentbench}.

\subsection{P1: deterministic execution aids, not privileged planning}

The language model should select goals and commitments, while deterministic tools handle brittle
mechanics already derivable from the local observation:

\begin{itemize}
  \item a local coordinate parser and short-path planner over visible tiles;
  \item semantic action validation against current affordances;
  \item a one-step ``face target'' helper;
  \item repeated-no-progress detection with obstacle evidence;
  \item an atomic action formatter that always emits one canonical action.
\end{itemize}

This is a tool-use treatment, not the canonical zero-shot policy.  It should be reported separately
because it changes the action scaffold.  The tool must receive no hidden state.  Compare invalid
actions, loops, navigation efficiency, coordination-site arrival, tokens, and final outcomes.  ReAct's
grounded reasoning/action cycle motivates this separation \cite{yao2023react}.

\subsection{P1: event-triggered and addressed communication}

Broadcasting every tick is expensive and can bury important changes.  Add a configurable policy
that sends only when an event occurs: new discovery, commitment state change, material inventory or
health threshold, failed action, request, or changed goal.  Allow direct addressing and local multicast;
retain a periodic heartbeat to detect silence.  Evaluate fixed byte budgets and induced delay/loss.

This treatment should answer whether communication quality, rather than volume, drives the paper's
ablation result.  Report value per delivered byte, redundant-message rate, missed critical events,
and team performance.  At larger populations, sparse routing is mandatory because all-to-all
broadcast grows quadratically.

\subsection{P1: failure detection and decentralized reassignment}

Introduce seeded crash-stop schedules distinct from in-world death.  A crashed agent receives no
model call, emits Noop, and communicates nothing.  Healthy agents must infer failure from missing
heartbeats or unmet commitments, cancel affected agreements, and negotiate takeover.  Measure
detection delay, abandoned obligations, reassignment latency, duplicated work, and recovered reward.

This provides a sharper test of adaptive collaboration than simply adding more agents.  It also
prepares the architecture for network partitions and timeouts without prematurely introducing
Byzantine behavior.

\subsection{P1: partner-generalization matrix}

Evaluate each policy with partners sampled across model, prompt, reasoning budget, message style,
and competence.  Use homogeneous controls, leave-one-partner-family-out tests, and intentionally
weak or silent partners.  Standardize the protocol schema across models while allowing free-text
payloads.  Report worst-partner, average-partner, and adaptation performance, not only homogeneous
mean.  This tests the coordination robustness emphasized by Overcooked and Melting Pot
\cite{carroll2019overcooked,leibo2021meltingpot} and avoids assuming that one strong model lifts a
mixed team.

\subsection{P2: skill learning and policy improvement}

Once diagnostics identify recurring failures, construct a versioned skill library from successful
traces: approach a shared target, stage a handover, execute a synchronous gather, transfer a resource,
recover from blocked navigation, and revive a teammate.  Skills should specify preconditions,
observable termination, and failure exits.  Retrieve only relevant skills and test held-out worlds.
Voyager supports the general value of executable skill accumulation \cite{wang2024voyager}, but
\dice{} should require multi-agent preconditions and commitments rather than copying a single-agent
library.

Subsequent training can proceed conservatively: supervised fine-tuning on validated decisions and
protocol messages; preference optimization against paired successful/failed coordination traces; then
offline or online RL using environment reward plus tightly scoped auxiliary rewards for valid
commitments.  Auxiliary shaping must not become the headline score, and held-out procedural seeds
and partner populations must remain untouched.

\section{Evaluation design}

\subsection{Preserve controls and use a staged funnel}

Each proposal starts with deterministic unit tests and a one-seed smoke test, then a three-seed
mechanism check, and only then the canonical 20 worlds on Easy, Medium, and Hard.  All material
comparisons use paired world seeds.  The published baseline remains frozen; new profiles include a
schema/version identifier and cannot write into baseline result directories.

The canonical metrics remain Team achievement percentage, coordination and normal achievement
percentages, action parse rate, reward decomposition, survival, usage, and latency.  Add the
opportunity funnel, commitment metrics, delivered bytes, and provider health.  Do not average the
three difficulties into one headline.  Report environment outcomes separately from failed attempts,
while retaining failed-attempt compute and error counts.

\subsection{Uncertainty and decision rules}

Twenty episodes are enough for a standard protocol but not necessarily a precise ranking.  Report
episode-level values, paired differences, bootstrap 95\% intervals, and robust aggregates such as the
interquartile mean.  Avoid declaring a win from point estimates whose intervals overlap materially.
This follows the warning that few-run RL comparisons can reverse under appropriate uncertainty
analysis \cite{agarwal2021precipice}.

Before running an experiment, register a primary outcome and guardrails.  For the commitment
protocol, a reasonable primary outcome is Hard synchronous-success per exposed opportunity.
Guardrails are Base\%, token count, delivered bytes, p95 tick latency, parse rate, and provider failure
rate.  Promotion requires a positive paired effect on the primary metric with no practically important
guardrail regression.  Exact thresholds should be fixed after measuring baseline variance, not after
seeing treatment results.

\subsection{Factorial ablations}

Avoid replacing the entire harness at once.  A minimal informative design is:

\begin{center}
\begin{tabular}{llll}
\toprule
Arm & Protocol & Structured memory & Deterministic tools \\
\midrule
A & free broadcast & no & no \\
B & commitment & no & no \\
C & free broadcast & yes & no \\
D & commitment & yes & no \\
E & commitment & yes & yes \\
\bottomrule
\end{tabular}
\end{center}

Run A--D first to estimate interaction between communication and memory.  Run E only if D shows a
coordination gain, because tools otherwise confound whether agreement or mechanical execution
improved.  Cross these arms with reasoning on/off only for models where the API implements a real
reasoning control.

\section{Serving and systems improvements}

Policy experiments are invalid when infrastructure failures selectively remove agents.  The observed
four-A100 setup illustrates three hazards: an Ollama server can be reachable yet lack the model;
automatic 256K context on an 80GB GPU can inflate the KV cache; and
\texttt{OLLAMA\_NUM\_PARALLEL} multiplies context memory \cite{ollamafaq,ollamacontext}.

The reliable single-episode topology is one server per GPU, one agent per server, explicit
\nolinkurl{CUDA_VISIBLE_DEVICES}, a shared read-only model store, explicit \texttt{num\_ctx}, and
\nolinkurl{OLLAMA_NUM_PARALLEL=1}.  A preflight must query \texttt{/api/tags}, preload the exact
digest, query \texttt{/api/ps}, and execute one bounded response on every endpoint.  Record model
digest, context allocation, GPU mapping, load time, prompt/eval tokens, time-to-first-token, generation
latency, queue time, and retry reason.

For multi-episode throughput, benchmark rather than assume.  Compare (i) one Ollama replica per GPU,
(ii) a batching-oriented vLLM deployment, and (iii) if supported, tensor-parallel serving.  vLLM's
PagedAttention was designed to reduce KV-cache waste and increase serving throughput
\cite{kwon2023vllm}; that does not guarantee a win for this exact Gemma build, so use identical
prompts and measure tokens/s, p50/p95 tick latency, failures, energy if available, and output
equivalence.  Separate leader latency from parallel worker-round latency.  Retry storms must be
bounded and counted; a 500 response should fail the preflight rather than consume an episode.

\section{Recommended implementation sequence}

\begin{longtable}{p{0.06\linewidth}p{0.30\linewidth}p{0.34\linewidth}p{0.14\linewidth}}
\toprule
Phase & Deliverable & Exit criterion & Priority \\
\midrule
\endhead
0 & Freeze baseline, add per-endpoint preflight and opportunity instrumentation. & Canonical trace unchanged; induced server faults classified before step 0. & \must \\
1 & Typed commitment envelopes and parser; no prompt/tool changes otherwise. & Unit tests plus paired smoke traces show correct lifecycle and no hidden-state access. & \must \\
2 & Structured patch-based private memory. & Lower duplication and stable commitment recall under long episodes. & \must \\
3 & A--D factorial on paired seeds. & Positive Hard opportunity-conversion evidence with uncertainty and guardrails. & \must \\
4 & Deterministic local execution aids and event-triggered routing. & Fewer mechanical failures and lower bytes without coordination regression. & \should \\
5 & Partner matrix and crash-stop schedules. & Performance characterized across unseen partners and failure fractions. & \should \\
6 & Skill library and training from validated trajectories. & Held-out seed/partner gains over strongest prompt-only arm. & \could \\
7 & Scale ladder at 8, 16, then 32 agents. & Bounded per-agent context, sparse communication, backpressure, and stable tick latency. & \could \\
\bottomrule
\end{longtable}

\section{Threats and non-recommendations}

First, changing prompts, memory, tools, and topology simultaneously prevents causal attribution.
Second, a bodyless leader can improve a pilot while weakening the claim of decentralized
coordination.  Third, more reasoning tokens may improve one model and cause another to omit the
action; reasoning controls are provider-specific.  Fourth, process metrics can be gamed if directly
rewarded, so environment outcomes remain primary.  Fifth, a single seed is a wiring or qualitative
study, not comparative evidence.  Sixth, training on canonical evaluation worlds or partners would
invalidate generalization claims.  Finally, higher GPU utilization is not itself an agent improvement;
systems changes must preserve request semantics and be reported separately.

The report does \emph{not} recommend silently replacing the standard baseline, treating message
length as communication quality, using hidden world state for navigation, assigning fixed roles from
the evaluator, or claiming large-scale collaboration from a crowded fixed map.  Those choices would
make scores easier to improve while weakening the scientific question.

\section{Conclusion}

\alem{} already identifies the right bottleneck: capable individual actors do not reliably form and
execute shared plans.  The most direct next experiment is therefore a decentralized commitment
layer paired with plan-oriented memory and opportunity-conditioned diagnostics.  This intervention
targets the paper's largest ablation effect, Gemma's observed memory weakness, and the synchronous-
hard failure chain without granting a planner privileged state.  Rigorous paired-seed evaluation,
unseen-partner tests, and fault-aware serving are equally important: without them, apparent gains can
be prompt variance, partner overfitting, or infrastructure selection effects.  The roadmap above turns
these ideas into separable, testable stages while preserving the published control.

\bibliographystyle{plain}
\bibliography{references}
\end{document}
